\documentclass[10pt,a4paper]{amsart}
\usepackage[margin=19mm]{geometry}
\usepackage{amsmath,amssymb,mathtools}
\usepackage[scr=rsfs,cal=euler]{mathalfa}
\usepackage{xfrac}
\usepackage[unicode,hidelinks,bookmarks=false]{hyperref}
\newcommand{\ie}{i.e.\ }
\newcommand{\cf}{cf.\ }
\newcommand{\resp}{respectively\ }
\newcommand{\Subsection}{\subsection}
\providecommand{\nobreakdash}{\nobreak\hbox{-}}
\setlength{\emergencystretch}{2em}
\multlinegap0pt
\usepackage{amssymb}
\newtheorem{lemm}{Lemma}
\newtheorem{cor}{Corollary}
\title{The Hausdorff dimension of the spectrum of the Thue-Morse Hamiltonian}
\author{Liu Qinghui, Lv Yiran}
\date{Source: arXiv:2609.12782v1 (CC BY 4.0)}
\begin{document}
\maketitle

\noindent\emph{Source and licence.} The Hausdorff dimension of the spectrum of the Thue-Morse Hamiltonian. Version arXiv:2609.12782v1 (CC BY 4.0), distributed by arXiv under the Creative Commons Attribution 4.0 International licence, http://creativecommons.org/licenses/by/4.0/, as recorded in the arXiv licence metadata for this exact version and shown on the abstract page https://arxiv.org/abs/2609.12782v1. Full licence notice: \texttt{licenses/arxiv-2609.12782-CC-BY-4.0.txt}.\par\smallskip

\noindent\textit{Editorial scope.} Counted unit: the article's cor 1924 (source0.tex lines 889--903). The article's complete original proof of it is delivered with it (source0.tex lines 906--940, one \texttt{proof} environment of 35 lines). No mathematics is written, repaired or re-derived: the statement and the proof are the article's own lines, in the article's own notation, and no retained line is substituted or re-flowed.  Retained uncounted as the source-local context the proof needs: lines 625--639; lines 876--884.  Nothing was omitted from any retained span, so the package carries no omission record.  No retained line contains a citation, so the wrapper carries no bibliography and no bibliography entry.  No retained line contains a figure environment, an image-inclusion command or drawing code, so no image is delivered and the package declares no asset dependency.  Editorial formatting only, in addition: an AMS \texttt{amsart} preamble that restates the article's own declarations and macros used by the retained lines, namely the environments \texttt{lemm}, \texttt{cor} on separate counters, together with the standard packages those lines need.  The retained environments print in this document as Lemm 1, Corollary 1 in order of appearance, whereas the article prints the same items with the numbering of its own full document.  The complete native source of the article as distributed in the arXiv e-print for this exact version is bundled unchanged as \texttt{sources/210152/source0.tex}.
\begin{lemm}
\label{lemma3}
For every integer $K\ge 1$, there exists $\delta_1>0$ such that, 
if $I_0$ is a band of level $n$  with $d(n,I_0)<\delta_1$, then  
\begin{equation*}
 I_0\cap\sigma_{n+1}\cap\cdots\cap\sigma_{n+K}
=I_1\sqcup\cdots\sqcup I_{2^K},
\end{equation*}
    where $I_1,\ldots,I_{2^K}$ are $2^K$ different intervals,
    and
\begin{equation*}
    x_{n+K}(I_j)\supset [-2,2\cos{\frac{\pi}{2^{K+1}}}]
\end{equation*}
    for  $ j=1,\ldots, 2^K$.
\end{lemm}

\begin{lemm}\label{discrete}
Let
$$Y_n=\{I_{n,j} : 0\le j<2^n \mbox{ and } S_n \varphi(t)\le-\log 80,\; \forall t\in I_{n,j} \}.$$
We have
$$\liminf_{n\to\infty}\frac{\sharp Y_n}{2^n}>\frac{1}{5}.$$
\end{lemm}
\paragraph{\textbf{Remark.}} Notice that $S_n\varphi(x)=S_n\varphi(1-x)$. 
Thus $Y_n$ is symmetric in the sense that
 $I_{n,j}\in Y_n$ whenever $I_{n,2^n-1-j}\in Y_n$. 

 \begin{cor}
\label{1924}
(1) There exists an integer $K_1>100$, such that for every $K>K_1$, 
\begin{equation}
\label{1/20}
\#\{I_{K,j}:\sup_{\theta\in I_{K,j}}R_K(\pi\theta)<1/20, 0\le j<2^K\}>2^K/6.
\end{equation}
(2) For $K>K_1$, there exists $\delta_3>0$, such that if $d(n,I_0)<\min\{\delta_1,\delta_3\}$
for a cut-off interval $I_0$ of level $n$, then there exists a symmetric family
\begin{equation*}
\mathcal{I}_{0,K}^{(2)}(I_0)\subset\{I\in\mathcal{I}_{0,K}^{(1)}(I_0):
d(n+K,I)<\frac{1}{10}d(n,I_0)\},
\end{equation*}
such that $\#\mathcal{I}_{0,K}^{(2)}(I_0)>2^K/6$.
\end{cor}

\begin{proof}
(1) Notice that 
\begin{equation*}
R_K(\pi\theta)=4\sin^2{(\pi\theta/2)}\mathrm{exp}(S_{K-1}\varphi(\theta))\le 4\mathrm{exp}(S_{K-1}\varphi(\theta)).
\end{equation*}
Then Lemma \ref{discrete} gives the proof.


(2) In the proof of Lemma \ref{lemma3},  
every $I_j\in \mathcal{I}_{0,K}^{(1)}(I_0)$ is a subset of $\xi(\pi I_{K,j-1})$.
     Furthermore, 
\begin{equation*}
    y_{n+K}= g(x_{n+K-1})\cdots g(x_{n+1})g(x_n)y_n.
    \end{equation*}
As in the proof of Lemma \ref{lemma3}, the function $g(x_{n+K-1})\cdots g(x_{n+1})g(x_n)$ is also
a polynomial in $x_n, y_n$, which we denote by $G(x_n,y_n)$. 
Thus there exists $\delta_3>0$ such that, for any $x\in [-2,2]$, 
for any $y$ with $|y|<\delta_{3}$, 
\begin{equation*}
|G(x,y)-G(x,0)|<1/20.
\end{equation*}
For $t\in I_0$, there is a $\theta\in[0,1]$ such that $x_n(t)=2\cos\pi\theta$.  By definition of $R_K$,
%Now take some $k\in\{0,2^K-1\}$ such that $P_k$ holds the condition \eqref{1/20}. We have
\begin{equation*}
G(x_n(t),0)=g(f^{K-1}(x_n(t)))\cdots g(f(x_n(t)))g(x_n(t))=R_{K}(\pi\theta).
\end{equation*}
If $R_K(\pi\theta)<1/20$, then
\begin{equation*}
|y_{n+K}(t)|=G(x_n(t),y_n(t))|y_n(t)|\le d(n,I_0)/10.
\end{equation*}
Define 
$$\mathcal{I}_{0,K}^{(2)}(I_0)=\{I_j\in\mathcal{I}_{0,K}^{(1)}(I_0) : 
I_{K,j-1}\in Y_n, j=1,\cdots, 2^K\}.$$
Hence, the fact that $\mathcal{I}_{0,K}^{(2)}(I_0)$ is symmetric and $\#\mathcal{I}_{0,K}^{(2)}(I_0)>2^K/6$ follows directly from remark of Lemma \ref{discrete} and  \eqref{1/20}.
\end{proof}

\end{document}
